\chapter{Objetivo general}


Desarrollar y validar un sistema predictivo híbrido que integre un modelo matemático basado en ecuaciones diferenciales ordinarias, sensores con tecnología IoT de bajo costo y algoritmos de aprendizaje automático, con el propósito de estimar en tiempo real las emisiones de gases de efecto invernadero (CO$_2$ y CH$_4$) durante el proceso de bioconversión de residuos orgánicos mediante larvas de \textit{Hermetia illucens}, contribuyendo a la optimización de la gestión ambiental y a la reducción del impacto climático en entornos agrícolas y agroindustriales.

\textbf{Objetivos Específicos}

\begin{enumerate}
    \item Diseñar y calibrar un sistema de sensores con IoT de bajo costo para el monitoreo en tiempo real de variables críticas durante el proceso de bioconversión con \textit{BSF}, como temperatura, humedad y concentraciones de CO$_2$ y CH$_4$.

    \item Modelar matemáticamente la generación de gases de efecto invernadero (GEI) en función del crecimiento larval, el consumo de sustrato y las condiciones ambientales, utilizando ecuaciones diferenciales ordinarias (EDO).

    \item Desarrollar e implementar algoritmos de aprendizaje automático que permitan la estimación y ajuste dinámico de los parámetros del modelo, utilizando los datos generados por los sensores en condiciones experimentales controladas.
    
    \item Validar el desempeño del sistema predictivo híbrido mediante la comparación con métodos convencionales de medición de GEI, evaluando su precisión, eficiencia y potencial de escalabilidad en escenarios agrícolas y agroindustriales.

\end{enumerate}
